\section{Discussion}
\markright{Discussion}

The experimental data shows that a Zener diode behaves like an ordinary diode in forward bias, where the current remains very small at low voltages and then rises rapidly once the forward threshold is reached. In reverse bias, the current remains almost zero until the breakdown region begins, after which the voltage across the diode remains nearly constant while the current increases significantly. From the measured values, the diode voltage stays close to about $5.4\,\text{V}$ in the breakdown region, which confirms the voltage regulation property of the Zener diode. The separate graphs for forward and reverse bias also show that the diode is suitable for maintaining a stable voltage under varying current conditions, although the resistance of the circuit and the applied supply must be chosen carefully to keep the diode within its safe operating range.



